\scanpage{240}\mspage{95}
Donc ces catégories, étant loc cofiltrantes (resp fbt),
sont en fait cofiltrantes (resp fbt cofiltrantes).

Prouvons c). Notons que pour $z$ dans $V$,
$\,{}_{z\backslash}U$ s'identifie à l'image inverse, par $\,{}_{z\backslash}(X_{/x})\unsure{-|\to}X_{/x}$,
de l'ouvert $U$ de $X_{/x}$, c'est donc un
\uline{ouvert} de $\,{}_{z\backslash}(X_{/x})$, non vide par construction
de $V$. Mais un ouvert non vide d'une
catégorie cofiltrante (resp. fbt. cofiltrante)
est encore cofiltrante (resp fbt. cofiltrante).
Cela signifie que $g:U\longrightarrow V$ est cofiltrant\deleted{\unsure{e}}
(resp fbt cofiltrant), a fortiori cofinal
(et $W$-coasphérique, pour tte $W$ satisfaisant
$\mathrm{Loc}(8)$). Cela achève la démonstration.

\uline{Remarque} \unsure{7.5}.\inserted{1)} Le même argument que dans
c) montre que si $f:X\unsure{-|\to}Y$ est
cofiltrant\unsure{e} (resp. fbt. cofiltrant), alors
pour \unsure{tout|tt} ouvert $U$ de $X$, si on
désigne par $V$ l'ouvert de $Y$ \deleted{\unsure{formé}} engendré
par $\unsure{f(U)}$, i.e. l'ouvert formé des $y$
de $Y$ tels que $\,{}_{y\backslash}X\ne\varnothing$, alors la
flèche induite $g:U\unsure{-|\to}V$ est cofiltrante
(resp. fbt cofiltrante), et \deleted{a fortiori}
est cofinal (et plus généralement, \inserted{si $f$ est cofiltrant, $g$ est}
$W$-coasphérique pour \deleted{\unsure{tout}}\inserted{tte} $W$ satisfaisant
\inserted{la condition} $\mathrm{Loc}(8)$).
% archive sheet number: 240
